% answers_short.tex —— 1995 年·全国卷·文科 · 简短答案速查
% 由 tools/make_short_answers.py 自 answers.tex 的 \gksolhead/\gkanswer 提取，
% 并与 1995/1995_数学真题_全国文_答案要点.md 逐题比对后生成，勿手改（改源 MD 后重跑）。

\begin{gknotice}
\textbf{使用说明}：本节只给最终答案便于快速核对，完整推导见解析版。来源页无答案；答案为题库独立推导（两名隔离求解员交叉核对一致），非卷内答案、非官方标准答案；第 16 题按对数内平方读法，第 24 题距离以底面半径表示。
\end{gknotice}

\gkpart{选择题}
\begin{center}\begin{tabular}{cccccccccccccccc}
\toprule
题号 & 1 & 2 & 3 & 4 & 5 & 6 & 7 & 8 & 9 & 10 & 11 & 12 & 13 & 14 & 15 \\
\midrule
答案 & B & D & C & B & D & C & A & C & A & A & B & D & D & C & A \\
\bottomrule
\end{tabular}\end{center}

\gkpart{填空题}
\begin{itemize}[leftmargin=2.6em]
  \item[\textbf{16.}] $x=3$。
  \item[\textbf{17.}] $\dfrac{7\sqrt3}{32}$。
  \item[\textbf{18.}] $\sqrt3$。
  \item[\textbf{19.}] $4$。
  \item[\textbf{20.}] $144$。
\end{itemize}

\gkpart{解答题}
\begin{itemize}[leftmargin=2.6em]
  \item[\textbf{21.}] $x=2$。
  \item[\textbf{22.}] $|z^2+z|=-2\cos\dfrac\theta2$（即 $2\left|\cos\dfrac\theta2\right|$）；辐角主值为 $\dfrac{3\theta}{2}-\pi$。
  \item[\textbf{23.}] 不等式对任意正整数 $n$ 成立（证明见解析）。
  \item[\textbf{24.}] （1）$AF\perp DB$；（2）点 $E$ 到截面 $ABCD$ 的距离等于圆柱底面半径。
  \item[\textbf{25.}] （1）$x=8+\dfrac{2\sqrt{50-t^2}-4t}{5}$，$t\in[0,\sqrt{10}]$；（2）政府补贴至少为每千克 $1$ 元。
  \item[\textbf{26.}] $(x-1)^2+\dfrac32y^2=1$（即 $2x^2+3y^2=4x$），$0<x\le2$，除去原点。
\end{itemize}
